\chapter{Chapter04 — DNF Fourier}
%%% generated by the blueprint pipeline from FABL.Chapter04.DNFFourier
\section{Overview}

Declarations of the Lean module \texttt{FABL.Chapter04.DNFFourier}.

\section{Declarations}

\begin{theorem}[Fourier coefficients of restricted functions]
\lean{FABL.fourierCoeff_extendedSignRestriction_liftFree}
\proofsource{Analysis-of-Boolean-Functions-by-Ryan-ODonnell}{
  pdf-ffbe9c7b3306-p076-b003,
  pdf-ffbe9c7b3306-p076-b004
}
\label{FABL.fourierCoeff_extendedSignRestriction_liftFree}
\leanok
\uses{FABL.FixedSignCube, FABL.FreeSignCube, FABL.combineSignCube, FABL.expect_freeProjection_signCubeSplitEquiv, FABL.extendedSignRestriction, FABL.fourierCoeff, FABL.freePart, FABL.freePart_eq_split, FABL.indexedMonomial, FABL.liftFreeFrequency, FABL.monomial, FABL.monomial_liftFreeFrequency_combine, FABL.restrictionFourierCoeff, FABL.signCubeSplitEquiv}
\difficulty{4}
Let $f$ be a real-valued function on the $n$-dimensional sign cube, let $J$ be a set of
coordinates, and fix signs on all coordinates outside $J$. Extend the restricted
function to the whole cube using those fixed signs outside $J$ and the input signs on
$J$. For any subset $S$ of $J$, the Fourier coefficient of this extension at the
corresponding coordinate set $S$ equals the Fourier coefficient of the restricted
function at $S$.
\end{theorem}

\begin{theorem}[Fourier coefficients under sign restriction]
\lean{FABL.fourierCoeff_extendedSignRestriction}
\proofsource{Analysis-of-Boolean-Functions-by-Ryan-ODonnell}{
  pdf-ffbe9c7b3306-p076-b003,
  pdf-ffbe9c7b3306-p076-b004
}
\label{FABL.fourierCoeff_extendedSignRestriction}
\leanok
\uses{FABL.FixedSignCube, FABL.ambientRestrictionFourierCoeff, FABL.extendedSignRestriction, FABL.fourierCoeff, FABL.fourierCoeff_extendedSignRestriction_liftFree, FABL.freeFrequencyPart, FABL.influence, FABL.influence_eq_sum_sq_fourierCoeff, FABL.influence_extendedSignRestriction_of_not_mem, FABL.liftFreeFrequency, FABL.liftFreeFrequency_freeFrequencyPart_of_subset, FABL.restrictionFourierCoeff}
\difficulty{5}
Let $f$ be a real-valued function on the $n$-dimensional sign cube, let $J$ be a set of
coordinates, and fix signs on all coordinates outside $J$. Extend the resulting
restricted function to the whole cube by using those fixed signs outside $J$ and the
input signs on $J$. For any coordinate set $S$, the Fourier coefficient of this
extension at $S$ equals the restricted Fourier coefficient at $S$ when $S\subseteq J$,
and equals zero otherwise.
\end{theorem}

\begin{theorem}[Fourier norm of an extended restriction]
\lean{FABL.sum_abs_ambientRestrictionFourierCoeff}
\label{FABL.sum_abs_ambientRestrictionFourierCoeff}
\leanok
\uses{FABL.FixedSignCube, FABL.ambientRestrictionFourierCoeff, FABL.extendedSignRestriction, FABL.fourierCoeff_extendedSignRestriction, FABL.fourierOneNorm}
\difficulty{3}
Let $n$ be a nonnegative integer, let $f$ be a real-valued function on the sign cube
$\{-1,1\}^n$, let $J\subseteq\{1,\ldots,n\}$, and fix signs $z$ on the coordinates
outside $J$. Extend the restriction of $f$ obtained by fixing those coordinates to a
function on the whole cube that ignores coordinates outside $J$. For each set of
coordinates $S$, take the Fourier coefficient of this extension when $S\subseteq J$,
and set it to zero otherwise; the sum of the absolute values of these coefficients over
all $S$ equals the Fourier one-norm of the extended function.
\end{theorem}

\begin{theorem}[Exercise4 11 restriction]
\lean{FABL.exercise4_11_restriction}
\statementsource{Analysis-of-Boolean-Functions-by-Ryan-ODonnell}{pdf-ffbe9c7b3306-p101-b006}
\label{FABL.exercise4_11_restriction}
\leanok
\uses{FABL.BooleanFunction, FABL.BooleanFunction.toReal, FABL.DecisionTree, FABL.FixedSignCube, FABL.F₂DecisionTree.Computes, FABL.F₂DecisionTree.computes_iff, FABL.F₂DecisionTree.depth, FABL.F₂DecisionTree.depth_mapOutputs, FABL.F₂DecisionTree.eval_mapOutputs, FABL.F₂DecisionTree.exists_computingTree_depth_eq_decisionTreeDepth, FABL.F₂DecisionTree.mapOutputs, FABL.F₂DecisionTree.spectralPNorm_one_le_infinityNorm_mul_two_pow_depth_of_computes, FABL.binaryCubeSignEquiv, FABL.binaryFunctionInfinityNorm, FABL.extendedSignRestriction, FABL.fourierOneNorm, FABL.fourierOneNorm_eq_spectralPNorm_one, FABL.restrictedBinaryFunction, FABL.restrictedDecisionTreeDepth, FABL.signValue, FABL.signValue_eq_neg_one_or_one, FABL.spectralPNorm}
\difficulty{5}
% dataset-pipeline: stub -- prose pending build_informal_statements.py
Statement pending regeneration.
\end{theorem}

\begin{theorem}[Absolute expectation of a random restriction]
\lean{FABL.abs_expectRandomRestriction_le_expect_abs}
\label{FABL.abs_expectRandomRestriction_le_expect_abs}
\leanok
\uses{FABL.FixedSignCube, FABL.deltaRandomSubsetWeight, FABL.deltaRandomSubsetWeight_nonneg, FABL.expectRandomRestriction}
\difficulty{4}
Let $n$ be a nonnegative integer, let $0\leq\delta\leq1$, and let $g$ assign a
real-valued function on each fixed sign cube indexed by a subset of the $n$ coordinates.
The absolute value of the expectation of $g$ over a random restriction is at most the
random-restriction expectation of the pointwise absolute value of $g$.
\end{theorem}

\begin{theorem}[Linearity of random restriction expectation]
\lean{FABL.sum_expectRandomRestriction}
\label{FABL.sum_expectRandomRestriction}
\leanok
\uses{FABL.FixedSignCube, FABL.expectRandomRestriction}
\difficulty{3}
Let $n$ be a nonnegative integer, let $\delta\in\mathbb R$, and let $\iota$ be a finite
index set. For each $a\in\iota$, let $g_a$ assign a real value to every choice of a
subset $J\subseteq\{0,\ldots,n-1\}$ and every sign assignment on the indices outside
$J$. The sum over $a$ of the random restriction expectations of the functions $g_a$
equals the random restriction expectation of their pointwise sum.
\end{theorem}

\begin{theorem}[Monotonicity of expected random restrictions]
\lean{FABL.expectRandomRestriction_mono}
\label{FABL.expectRandomRestriction_mono}
\leanok
\uses{FABL.FixedSignCube, FABL.deltaRandomSubsetWeight_nonneg, FABL.expectRandomRestriction}
\difficulty{3}
Let $n$ be a nonnegative integer and let $0\leq\delta\leq1$. For each subset
$J\subseteq\{1,\ldots,n\}$, let a fixed sign assignment on the indices outside $J$
determine a real-valued function, and suppose $g(J,z)\leq h(J,z)$ for every such $J$ and
assignment $z$. Then the expected random restriction of $g$, with each subset weighted
by $\delta^{|J|}(1-\delta)^{n-|J|}$ and averaged over the remaining signs, is at most
the corresponding expected random restriction of $h$.
\end{theorem}

\begin{theorem}[Fourier mass under random restrictions]
\lean{FABL.lemma4_23}
\proofsource{Analysis-of-Boolean-Functions-by-Ryan-ODonnell}{pdf-ffbe9c7b3306-p102-b001}
\label{FABL.lemma4_23}
\leanok
\uses{FABL.BooleanFunction, FABL.BooleanFunction.toReal, FABL.abs_expectRandomRestriction_le_expect_abs, FABL.ambientRestrictionFourierCoeff, FABL.exercise4_11_restriction, FABL.expectRandomRestriction, FABL.expectRandomRestriction_mono, FABL.expect_fourierCoeff_randomRestriction, FABL.extendedSignRestriction, FABL.fourierCoeff, FABL.fourierOneNorm, FABL.restrictedDecisionTreeDepth, FABL.sum_abs_ambientRestrictionFourierCoeff, FABL.sum_expectRandomRestriction}
\difficulty{4}
Let $f:\{-1,1\}^n\to\{-1,1\}$ and $0\leq\delta\leq1$. For each $S\subseteq[n]$, let
$\widehat f(S)$ be the uniform Fourier coefficient of $f$ at the character
$x\mapsto\prod_{i\in S}x_i$. Then $\sum_{S\subseteq[n]}\delta^{|S|}|\widehat f(S)|$ is
at most the expected value of $2^D$ under a random restriction that leaves each
coordinate free independently with probability $\delta$ and fixes every other coordinate
to an independent uniform sign, where $D$ is the decision-tree depth of the restricted
function.
\end{theorem}

\begin{definition}[Delta random intersection tail]
\lean{FABL.deltaRandomIntersectionTail}
\statementsource{Analysis-of-Boolean-Functions-by-Ryan-ODonnell}{pdf-ffbe9c7b3306-p101-b003}
\label{FABL.deltaRandomIntersectionTail}
\leanok
\uses{FABL.deltaRandomSubsetWeight}
For a natural number $n$, a real parameter $\delta$, a subset
$U\subseteq\mathrm{Fin}(n)$, and a natural number $k$, define the delta random
intersection tail to be $\sum_{J\subseteq\mathrm{Fin}(n)}
w_{n,\delta}(J)\,\mathbf{1}_{\{k\le |U\cap J|\}}$, where $w_{n,\delta}(J)$ denotes the
delta random subset weight of $J$.
\end{definition}

\begin{theorem}[Intersection cardinality as an indicator sum]
\lean{FABL.card_inter_eq_sum_memIndicator}
\label{FABL.card_inter_eq_sum_memIndicator}
\leanok
\difficulty{4}
Let $n$ be a natural number and let $U,J$ be finite subsets of $\{0,\ldots,n-1\}$.
The cardinality of their intersection, viewed as a real number, is $\sum_{i\in
U}\mathbf{1}_{J}(i)$, where $\mathbf{1}_{J}(i)$ is $1$ if $i\in J$ and $0$ otherwise.
\end{theorem}

\begin{theorem}[Weight of subsets containing a pair]
\lean{FABL.sum_deltaRandomSubsetWeight_pair_mem}
\label{FABL.sum_deltaRandomSubsetWeight_pair_mem}
\leanok
\uses{FABL.deltaRandomSubsetWeight, FABL.sum_deltaRandomSubsetWeight_supset}
\difficulty{4}
Let $n$ be a natural number, let $\delta\in\mathbb{R}$, and let $i,j$ be elements of the
$n$-element set $\{0,\ldots,n-1\}$. Give each subset $J$ of this set weight
$w(J)=\delta^{|J|}(1-\delta)^{n-|J|}$. The sum of the weights of subsets containing both $i$
and $j$ is $\delta$ if $i=j$, and $\delta^2$ otherwise.
\end{theorem}

\begin{theorem}[Expected subset intersection size]
\lean{FABL.sum_deltaRandomSubsetWeight_mul_card_inter}
\label{FABL.sum_deltaRandomSubsetWeight_mul_card_inter}
\leanok
\uses{FABL.card_inter_eq_sum_memIndicator, FABL.deltaRandomSubsetWeight, FABL.sum_deltaRandomSubsetWeight_mem}
\difficulty{4}
Let $n$ be a nonnegative integer, let $\delta\in\mathbb R$, and let
$U\subseteq\{0,\ldots,n-1\}$. Weight each subset $J$ of the coordinate set by
$\delta^{|J|}(1-\delta)^{n-|J|}$. Then $\sum_J \delta^{|J|}(1-\delta)^{n-|J|}|U\cap
J|=\delta|U|$.
\end{theorem}

\begin{theorem}[Intersection Cardinality as a Pair Count]
\lean{FABL.sq_card_inter_eq_sum_pairIndicator}
\label{FABL.sq_card_inter_eq_sum_pairIndicator}
\leanok
\uses{FABL.card_inter_eq_sum_memIndicator}
\difficulty{4}
For subsets $U,J\subseteq\{1,\ldots,n\}$, the square of the size of their intersection
equals the number of ordered pairs $(i,j)\in U\times U$ for which both $i$ and $j$
belong to $J$: $|U\cap J|^2=\sum_{i\in U}\sum_{j\in U}\mathbf{1}_{i\in J\,\wedge\,j\in
J}$.
\end{theorem}

\begin{theorem}[Sum of pairwise diagonal weights]
\lean{FABL.sum_pairDiagonalWeight}
\label{FABL.sum_pairDiagonalWeight}
\leanok
\difficulty{4}
Let $n$ be a natural number, let $\delta\in\mathbb R$, and let $U$ be a finite subset of
the $n$ coordinates. The sum over ordered pairs $(i,j)\in U\times U$, assigning weight
$\delta$ when $i=j$ and $\delta^2$ otherwise, equals
$\delta^2|U|^2+\delta(1-\delta)|U|$.
\end{theorem}

\begin{theorem}[Sum delta random subset weight mul sq card inter]
\lean{FABL.sum_deltaRandomSubsetWeight_mul_sq_card_inter}
\statementsource{Analysis-of-Boolean-Functions-by-Ryan-ODonnell}{pdf-ffbe9c7b3306-p101-b003}
\label{FABL.sum_deltaRandomSubsetWeight_mul_sq_card_inter}
\leanok
\uses{FABL.deltaRandomSubsetWeight, FABL.sq_card_inter_eq_sum_pairIndicator, FABL.sum_deltaRandomSubsetWeight_pair_mem, FABL.sum_pairDiagonalWeight}
\difficulty{4}
% dataset-pipeline: stub -- prose pending build_informal_statements.py
Statement pending regeneration.
\end{theorem}

\begin{theorem}[Weighted intersection-size variance identity]
\lean{FABL.sum_deltaRandomSubsetWeight_mul_sq_card_inter_sub_mean}
\label{FABL.sum_deltaRandomSubsetWeight_mul_sq_card_inter_sub_mean}
\leanok
\uses{FABL.deltaRandomSubsetWeight, FABL.sum_deltaRandomSubsetWeight, FABL.sum_deltaRandomSubsetWeight_mul_card_inter, FABL.sum_deltaRandomSubsetWeight_mul_sq_card_inter}
\difficulty{4}
Let $n$ be a natural number, let $\delta\in\mathbb{R}$, and let
$U\subseteq\{0,\ldots,n-1\}$. For each subset $J\subseteq\{0,\ldots,n-1\}$, set
$w_\delta(J)=\delta^{|J|}(1-\delta)^{n-|J|}$. Then \[\sum_{J\subseteq\{0,\ldots,n-1\}}
w_\delta(J)\bigl(|U\cap J|-\delta|U|\bigr)^2=\delta(1-\delta)|U|.\]
\end{theorem}

\begin{theorem}[One third le delta random intersection tail]
\lean{FABL.one_third_le_deltaRandomIntersectionTail}
\statementsource{Analysis-of-Boolean-Functions-by-Ryan-ODonnell}{pdf-ffbe9c7b3306-p101-b003}
\label{FABL.one_third_le_deltaRandomIntersectionTail}
\leanok
\uses{FABL.deltaRandomIntersectionTail, FABL.deltaRandomSubsetWeight, FABL.deltaRandomSubsetWeight_nonneg, FABL.sum_deltaRandomSubsetWeight, FABL.sum_deltaRandomSubsetWeight_mul_sq_card_inter_sub_mean}
\difficulty{6}
% dataset-pipeline: stub -- prose pending build_informal_statements.py
Statement pending regeneration.
\end{theorem}

\begin{theorem}[Nonnegativity of the intersection tail]
\lean{FABL.deltaRandomIntersectionTail_nonneg}
\label{FABL.deltaRandomIntersectionTail_nonneg}
\leanok
\uses{FABL.deltaRandomIntersectionTail, FABL.deltaRandomSubsetWeight_nonneg}
\difficulty{4}
Let $n\in\mathbb{N}$, let $\delta\in\mathbb{R}$ satisfy $0\le\delta\le1$, let
$U\subseteq\{0,\ldots,n-1\}$, and let $k\in\mathbb{N}$. The weighted sum of the
indicators that a random subset's intersection with $U$ has size at least $k$ is
nonnegative.
\end{theorem}

\begin{theorem}[Subset-sum formula for intersection tails]
\lean{FABL.deltaRandomIntersectionTail_eq_sum}
\label{FABL.deltaRandomIntersectionTail_eq_sum}
\leanok
\uses{FABL.deltaRandomIntersectionTail, FABL.deltaRandomSubsetWeight, FABL.sum_deltaRandomSubsetWeight_inter_eq}
\difficulty{5}
Let $n,k\in\mathbb{N}$, let $\delta\in\mathbb{R}$, and let $U\subseteq\{0,\ldots,n-1\}$.
The weighted sum over subsets $J\subseteq\{0,\ldots,n-1\}$ whose intersection with $U$
has size at least $k$ equals $\sum_{S\subseteq\{0,\ldots,n-1\}}\mathbf{1}_{\{k\le
|S|\}}\mathbf{1}_{\{S\subseteq U\}}\delta^{|S|}(1-\delta)^{|U\setminus S|}$, where
$\mathbf{1}$ denotes the indicator, with value $1$ when its condition holds and $0$
otherwise.
\end{theorem}

\begin{definition}[Restricted Fourier weight above a degree threshold]
\lean{FABL.restrictedFourierWeightAtLeast}
\statementsource{Analysis-of-Boolean-Functions-by-Ryan-ODonnell}{pdf-ffbe9c7b3306-p101-b003}
\label{FABL.restrictedFourierWeightAtLeast}
\leanok
\uses{FABL.BooleanFunction, FABL.BooleanFunction.toReal, FABL.FixedSignCube, FABL.ambientRestrictionFourierCoeff}
Let $n$ be a natural number, let $f$ be a Boolean function on the $n$-dimensional sign
cube, let $k$ be a natural number, let $J\subseteq\mathrm{Fin}(n)$, and fix a sign
assignment $z$ on the indices outside $J$. Define the restricted Fourier weight at
degrees at least $k$ to be $\sum_{S\subseteq\mathrm{Fin}(n)}\mathbf{1}_{\{k\le
|S|\}}\,\widehat f_{J,z}(S)^2$, where $\widehat f_{J,z}(S)$ denotes the ambient
restriction Fourier coefficient of $f$ at $J,S,z$.
\end{definition}

\begin{theorem}[Degree under an extended sign restriction]
\lean{FABL.fourierDegree_extendedSignRestriction_le_restrictedDecisionTreeDepth}
\proofsource{Analysis-of-Boolean-Functions-by-Ryan-ODonnell}{
  pdf-ffbe9c7b3306-p074-b005,
  pdf-ffbe9c7b3306-p101-b003
}
\label{FABL.fourierDegree_extendedSignRestriction_le_restrictedDecisionTreeDepth}
\leanok
\uses{FABL.BooleanFunction, FABL.BooleanFunction.toReal, FABL.DecisionTree, FABL.FixedSignCube, FABL.F₂DecisionTree.Computes, FABL.F₂DecisionTree.computes_iff, FABL.F₂DecisionTree.depth_mapOutputs, FABL.F₂DecisionTree.eval_mapOutputs, FABL.F₂DecisionTree.exists_computingTree_depth_eq_decisionTreeDepth, FABL.F₂DecisionTree.mapOutputs, FABL.binaryCubeSignEquiv, FABL.extendedSignRestriction, FABL.fourierDegree, FABL.fourierDegree_toReal_le_depth_of_decisionTree, FABL.restrictedBinaryFunction, FABL.restrictedDecisionTreeDepth, FABL.signValue}
\difficulty{4}
Let $n$ be a natural number, let $f\colon\{-1,1\}^n\to\{-1,1\}$, and choose a set $J$ of
coordinates together with fixed signs on the coordinates outside $J$. Extend this
restriction to a function on the whole sign cube by evaluating $f$ at the point whose
coordinates in $J$ come from the input and whose other coordinates are the chosen fixed
signs. Its real-valued Fourier degree is at most the depth of a decision tree computing
the corresponding restricted function on the binary cube, where binary inputs are
identified with signs.
\end{theorem}

\begin{theorem}[Restricted fourier weight at least le switching failure indicator]
\lean{FABL.restrictedFourierWeightAtLeast_le_switchingFailureIndicator}
\statementsource{Analysis-of-Boolean-Functions-by-Ryan-ODonnell}{pdf-ffbe9c7b3306-p101-b003}
\label{FABL.restrictedFourierWeightAtLeast_le_switchingFailureIndicator}
\leanok
\uses{FABL.BooleanFunction, FABL.BooleanFunction.toReal, FABL.FixedSignCube, FABL.ambientRestrictionFourierCoeff, FABL.extendedSignRestriction, FABL.fourierCoeff, FABL.fourierCoeff_extendedSignRestriction, FABL.fourierDegree_extendedSignRestriction_le_restrictedDecisionTreeDepth, FABL.fourierDegree_le_iff, FABL.restrictedDecisionTreeDepth, FABL.restrictedFourierWeightAtLeast, FABL.sum_sq_fourierCoeff_eq_one, FABL.switchingFailureIndicator}
\difficulty{5}
% dataset-pipeline: stub -- prose pending build_informal_statements.py
Statement pending regeneration.
\end{theorem}

\begin{theorem}[Expect restricted fourier weight at least le switching failure probability]
\lean{FABL.expect_restrictedFourierWeightAtLeast_le_switchingFailureProbability}
\statementsource{Analysis-of-Boolean-Functions-by-Ryan-ODonnell}{pdf-ffbe9c7b3306-p101-b003}
\label{FABL.expect_restrictedFourierWeightAtLeast_le_switchingFailureProbability}
\leanok
\uses{FABL.BooleanFunction, FABL.expectRandomRestriction, FABL.expectRandomRestriction_mono, FABL.restrictedFourierWeightAtLeast, FABL.restrictedFourierWeightAtLeast_le_switchingFailureIndicator, FABL.switchingFailureProbability}
\difficulty{3}
% dataset-pipeline: stub -- prose pending build_informal_statements.py
Statement pending regeneration.
\end{theorem}

\begin{theorem}[Expect restricted fourier weight at least eq sum intersection tail]
\lean{FABL.expect_restrictedFourierWeightAtLeast_eq_sum_intersectionTail}
\statementsource{Analysis-of-Boolean-Functions-by-Ryan-ODonnell}{
  pdf-ffbe9c7b3306-p099-b001,
  pdf-ffbe9c7b3306-p101-b003
}
\label{FABL.expect_restrictedFourierWeightAtLeast_eq_sum_intersectionTail}
\leanok
\uses{FABL.BooleanFunction, FABL.BooleanFunction.toReal, FABL.ambientRestrictionFourierCoeff, FABL.deltaRandomIntersectionTail, FABL.deltaRandomIntersectionTail_eq_sum, FABL.expectRandomRestriction, FABL.expect_sq_fourierCoeff_randomRestriction, FABL.fourierCoeff, FABL.restrictedFourierWeightAtLeast, FABL.sum_expectRandomRestriction}
\difficulty{5}
% dataset-pipeline: stub -- prose pending build_informal_statements.py
Statement pending regeneration.
\end{theorem}

\begin{theorem}[Lemma4 21]
\lean{FABL.lemma4_21}
\statementsource{Analysis-of-Boolean-Functions-by-Ryan-ODonnell}{
  pdf-ffbe9c7b3306-p101-b002,
  pdf-ffbe9c7b3306-p101-b003
}
\label{FABL.lemma4_21}
\leanok
\uses{FABL.BooleanFunction, FABL.BooleanFunction.toReal, FABL.IsFourierSpectrumConcentratedUpTo, FABL.deltaRandomIntersectionTail, FABL.deltaRandomIntersectionTail_nonneg, FABL.expect_restrictedFourierWeightAtLeast_eq_sum_intersectionTail, FABL.expect_restrictedFourierWeightAtLeast_le_switchingFailureProbability, FABL.fourierCoeff, FABL.fourierWeightAboveReal, FABL.one_third_le_deltaRandomIntersectionTail, FABL.switchingFailureProbability}
\difficulty{4}
% dataset-pipeline: stub -- prose pending build_informal_statements.py
Statement pending regeneration.
\end{theorem}

\begin{theorem}[Theorem4 22 of switching failure probability le]
\lean{FABL.theorem4_22_of_switchingFailureProbability_le}
\statementsource{Analysis-of-Boolean-Functions-by-Ryan-ODonnell}{
  pdf-ffbe9c7b3306-p101-b002,
  pdf-ffbe9c7b3306-p101-b003,
  pdf-ffbe9c7b3306-p101-b005
}
\label{FABL.theorem4_22_of_switchingFailureProbability_le}
\leanok
\uses{FABL.BooleanFunction, FABL.BooleanFunction.toReal, FABL.IsFourierSpectrumConcentratedUpTo, FABL.IsFourierSpectrumConcentratedUpTo.mono_error, FABL.lemma4_21, FABL.switchingFailureProbability}
\difficulty{3}
% dataset-pipeline: stub -- prose pending build_informal_statements.py
Statement pending regeneration.
\end{theorem}

\begin{theorem}[Restricted decision tree depth le dimension]
\lean{FABL.restrictedDecisionTreeDepth_le_dimension}
\statementsource{Analysis-of-Boolean-Functions-by-Ryan-ODonnell}{pdf-ffbe9c7b3306-p074-b002}
\label{FABL.restrictedDecisionTreeDepth_le_dimension}
\leanok
\uses{FABL.BooleanFunction, FABL.FixedSignCube, FABL.F₂DecisionTree.depth_le_dimension, FABL.F₂DecisionTree.exists_computingTree_depth_eq_decisionTreeDepth, FABL.restrictedBinaryFunction, FABL.restrictedDecisionTreeDepth}
\difficulty{4}
% dataset-pipeline: stub -- prose pending build_informal_statements.py
Statement pending regeneration.
\end{theorem}

\begin{theorem}[Expectation of a constant under random restriction]
\lean{FABL.expectRandomRestriction_const}
\label{FABL.expectRandomRestriction_const}
\leanok
\uses{FABL.expectRandomRestriction, FABL.sum_deltaRandomSubsetWeight}
\difficulty{3}
For every nonnegative integer $n$ and all real numbers $\delta$ and $c$, the random
restriction expectation of the constant function with value $c$ is $c$.
\end{theorem}

\begin{theorem}[Additivity of random restriction expectation]
\lean{FABL.expectRandomRestriction_add}
\label{FABL.expectRandomRestriction_add}
\leanok
\uses{FABL.FixedSignCube, FABL.expectRandomRestriction}
\difficulty{3}
Let $n$ be a nonnegative integer and let $\delta\in\mathbb R$. For any real-valued
functions $g$ and $h$ that assign a value to each choice of a free-coordinate subset
$J\subseteq\{1,\ldots,n\}$ and signs on its complementary coordinates, the random
restriction expectation of their pointwise sum is the sum of their random restriction
expectations.
\end{theorem}

\begin{theorem}[Homogeneity of random restriction expectation]
\lean{FABL.expectRandomRestriction_const_mul}
\label{FABL.expectRandomRestriction_const_mul}
\leanok
\uses{FABL.FixedSignCube, FABL.deltaRandomSubsetWeight, FABL.expectRandomRestriction}
\difficulty{4}
Let $n$ be a nonnegative integer, let $\delta,c\in\mathbb R$, and let $g$ be any
real-valued function of a free-coordinate subset of $\{1,\ldots,n\}$ and signs on its
complement. The random restriction expectation of $c$ times $g$ is $c$ times the random
restriction expectation of $g$.
\end{theorem}

\begin{theorem}[Linearity of expected random restrictions]
\lean{FABL.expectRandomRestriction_sum}
\label{FABL.expectRandomRestriction_sum}
\leanok
\uses{FABL.FixedSignCube, FABL.expectRandomRestriction, FABL.expectRandomRestriction_add, FABL.expectRandomRestriction_const}
\difficulty{4}
Let $n$ be a nonnegative integer, let $\delta$ be real, let $\iota$ be a type, and let
$s$ be a finite subset of $\iota$. For each $a\in\iota$, let $g_a(J,z)$ be a real-valued
function of a subset $J\subseteq\{1,\ldots,n\}$ and a sign assignment $z$ on its
complement. The expected random restriction of the pointwise sum $\sum_{a\in s}g_a$
equals the sum over $a\in s$ of the expected random restrictions of $g_a$.
\end{theorem}

\begin{theorem}[Finite geometric sum for powers of two]
\lean{FABL.two_pow_eq_one_add_sum_range}
\label{FABL.two_pow_eq_one_add_sum_range}
\leanok
\difficulty{3}
For every natural number $d$, $2^d=1+\sum_{j=0}^{d-1}2^j$, with the natural-number
powers interpreted as real numbers.
\end{theorem}

\begin{theorem}[Power of two as a truncated indicator sum]
\lean{FABL.two_pow_eq_one_add_sum_tailIndicators}
\label{FABL.two_pow_eq_one_add_sum_tailIndicators}
\leanok
\uses{FABL.two_pow_eq_one_add_sum_range}
\difficulty{4}
Let $d$ and $N$ be natural numbers with $d\le N$. Then
$2^d=1+\sum_{j=0}^{N-1}2^j\,\boldsymbol{1}_{\{j+1\le d\	ext	ext{}}$, where each
natural-number power is interpreted as a real number and the indicator is $1$ when
$j+1\le d$ and $0$ otherwise.
\end{theorem}

\begin{theorem}[Decision tree depth and tail indicators]
\lean{FABL.two_pow_restrictedDecisionTreeDepth_eq_tailIndicators}
\label{FABL.two_pow_restrictedDecisionTreeDepth_eq_tailIndicators}
\leanok
\uses{FABL.BooleanFunction, FABL.FixedSignCube, FABL.restrictedDecisionTreeDepth, FABL.restrictedDecisionTreeDepth_le_dimension, FABL.switchingFailureIndicator, FABL.two_pow_eq_one_add_sum_tailIndicators}
\difficulty{3}
Let $n$ be a nonnegative integer, let $f:\{\pm1\}^n\to\{\pm1\}$ be a Boolean function,
let $J\subseteq\{1,\ldots,n\}$ be a set of coordinates, and fix signs on the coordinates
outside $J$. Write $D$ for the depth of the decision tree computing the resulting
restricted function, and for each $k$ let $I_k$ be $1$ if $k\le D$ and $0$ otherwise.
Then $2^D=1+\sum_{j=0}^{n-1}2^j I_{j+1}$.
\end{theorem}

\begin{theorem}[Depth moment and failure probabilities]
\lean{FABL.expect_two_pow_restrictedDecisionTreeDepth_eq_failureSum}
\label{FABL.expect_two_pow_restrictedDecisionTreeDepth_eq_failureSum}
\leanok
\uses{FABL.BooleanFunction, FABL.expectRandomRestriction, FABL.expectRandomRestriction_add, FABL.expectRandomRestriction_const, FABL.expectRandomRestriction_const_mul, FABL.expectRandomRestriction_sum, FABL.restrictedDecisionTreeDepth, FABL.switchingFailureIndicator, FABL.switchingFailureProbability, FABL.two_pow_restrictedDecisionTreeDepth_eq_tailIndicators}
\difficulty{4}
Let $n$ be a nonnegative integer, let $f$ be a sign-valued function on the
$n$-dimensional sign cube, and let $\delta$ be any real number. For each subset $J$ of
coordinates and each assignment $z$ of signs outside $J$, let $d(J,z)$ be the minimum
depth of a binary decision tree computing the function obtained by fixing those outside
coordinates; define $q_k$ to be the random-restriction average of the indicator of $k\le
d(J,z)$, using weight $\delta^{|J|}(1-\delta)^{n-|J|}$ for $J$ and a uniform average
over $z$. Then the random-restriction average of $2^{d(J,z)}$ is $1+\sum_{j=0}^{n-1}2^j
q_{j+1}$.
\end{theorem}

\begin{theorem}[A power identity]
\lean{FABL.two_pow_mul_quarter_pow_succ}
\label{FABL.two_pow_mul_quarter_pow_succ}
\leanok
\difficulty{3}
For every nonnegative integer $j$, $2^j(1/4)^{j+1}=(1/4)(1/2)^j$.
\end{theorem}

\begin{theorem}[Expect two pow restricted decision tree depth le two of failure le quarter]
\lean{FABL.expect_two_pow_restrictedDecisionTreeDepth_le_two_of_failure_le_quarter}
\statementsource{Analysis-of-Boolean-Functions-by-Ryan-ODonnell}{
  pdf-ffbe9c7b3306-p100-b005,
  pdf-ffbe9c7b3306-p102-b003
}
\label{FABL.expect_two_pow_restrictedDecisionTreeDepth_le_two_of_failure_le_quarter}
\leanok
\uses{FABL.BooleanFunction, FABL.expectRandomRestriction, FABL.expect_two_pow_restrictedDecisionTreeDepth_eq_failureSum, FABL.restrictedDecisionTreeDepth, FABL.switchingFailureProbability, FABL.two_pow_mul_quarter_pow_succ}
\difficulty{4}
% dataset-pipeline: stub -- prose pending build_informal_statements.py
Statement pending regeneration.
\end{theorem}

\begin{theorem}[Low degree Fourier coefficient bound]
\lean{FABL.lowDegreeFourierOneNorm_le_two_mul_inv_pow_of_expected_two_pow_le}
\proofsource{Analysis-of-Boolean-Functions-by-Ryan-ODonnell}{
  pdf-ffbe9c7b3306-p102-b001,
  pdf-ffbe9c7b3306-p102-b003
}
\label{FABL.lowDegreeFourierOneNorm_le_two_mul_inv_pow_of_expected_two_pow_le}
\leanok
\uses{FABL.BooleanFunction, FABL.BooleanFunction.toReal, FABL.expectRandomRestriction, FABL.fourierCoeff, FABL.lemma4_23, FABL.restrictedDecisionTreeDepth}
\difficulty{4}
Let $n$ and $k$ be nonnegative integers, let $f$ be a Boolean function on the
$n$-dimensional sign cube, and let $0<\delta\leq 1$. If the expected value, over the
random restriction, of $2$ raised to the restricted function’s decision tree depth is at
most $2$, then the sum of the absolute values of the Fourier coefficients of $f$ over
all subsets of size at most $k$ is at most $2\delta^{-k}$.
\end{theorem}

\begin{theorem}[Low-degree Fourier bound from switching failure]
\lean{FABL.theorem4_24_of_switchingFailureProbability_le_quarter}
\proofsource{Analysis-of-Boolean-Functions-by-Ryan-ODonnell}{
  pdf-ffbe9c7b3306-p102-b002,
  pdf-ffbe9c7b3306-p102-b003
}
\label{FABL.theorem4_24_of_switchingFailureProbability_le_quarter}
\leanok
\uses{FABL.BooleanFunction, FABL.BooleanFunction.toReal, FABL.expectRandomRestriction, FABL.expect_two_pow_restrictedDecisionTreeDepth_le_two_of_failure_le_quarter, FABL.fourierCoeff, FABL.lowDegreeFourierOneNorm_le_two_mul_inv_pow_of_expected_two_pow_le, FABL.restrictedDecisionTreeDepth, FABL.switchingFailureProbability}
\difficulty{4}
Let $n$ be a nonnegative integer and let $f$ be a sign-valued Boolean function on the
$n$-dimensional sign cube. Suppose $w$ is a positive integer and, for every nonnegative
integer $d$, the probability that a random restriction with retention parameter
$1/(20w)$ leaves decision-tree depth at least $d$ is at most $(1/4)^d$. Then for every
nonnegative integer $k$, the sum of the absolute values of the Fourier coefficients of
the real-valued version of $f$ over all coordinate subsets of size at most $k$ is at
most $2(20w)^k$.
\end{theorem}

\begin{definition}[DNF switching depth]
\lean{FABL.dnfSwitchingDepth}
\statementsource{Analysis-of-Boolean-Functions-by-Ryan-ODonnell}{pdf-ffbe9c7b3306-p100-b005}
\label{FABL.dnfSwitchingDepth}
\leanok
For a real number $\varepsilon$, define the DNF switching depth to be the natural-number
ceiling of $\log_2(3/\varepsilon)$.
\end{definition}

\begin{theorem}[Switching-Depth Power Bound]
\lean{FABL.half_pow_dnfSwitchingDepth_le}
\label{FABL.half_pow_dnfSwitchingDepth_le}
\leanok
\uses{FABL.dnfSwitchingDepth}
\difficulty{4}
For every real number $\varepsilon>0$, let $d=\lceil\log_2(3/\varepsilon)\rceil$. Then
$(1/2)^d\le \varepsilon/3$.
\end{theorem}

\begin{theorem}[Positivity of the switching depth]
\lean{FABL.dnfSwitchingDepth_pos}
\label{FABL.dnfSwitchingDepth_pos}
\leanok
\uses{FABL.dnfSwitchingDepth}
\difficulty{3}
For a real number $\varepsilon$ with $0<\varepsilon\le 1$, the switching depth
$\lceil\log_2(3/\varepsilon)\rceil$ is positive.
\end{theorem}

\begin{theorem}[Fourier concentration from switching failure bounds]
\lean{FABL.theorem4_22_of_switchingFailureProbability_le_half_pow}
\proofsource{Analysis-of-Boolean-Functions-by-Ryan-ODonnell}{
  pdf-ffbe9c7b3306-p101-b004,
  pdf-ffbe9c7b3306-p101-b005,
  pdf-ffbe9c7b3306-p101-b002,
  pdf-ffbe9c7b3306-p101-b003
}
\label{FABL.theorem4_22_of_switchingFailureProbability_le_half_pow}
\leanok
\uses{FABL.BooleanFunction, FABL.BooleanFunction.toReal, FABL.IsFourierSpectrumConcentratedUpTo, FABL.dnfSwitchingDepth, FABL.dnfSwitchingDepth_pos, FABL.half_pow_dnfSwitchingDepth_le, FABL.switchingFailureProbability, FABL.theorem4_22_of_switchingFailureProbability_le}
\difficulty{4}
Let $n$ be a nonnegative integer and let $f$ be a sign-valued Boolean function on the
$n$-dimensional sign cube. Suppose $f$ has a DNF representation of width at most the
positive integer $w$, and let $0<\varepsilon\le1$. Assume that for every nonnegative
integer $d$, the probability that a random restriction with retention parameter
$1/(10w)$ leaves decision-tree depth at least $d$ is at most $(1/2)^d$. Then the sum of
squared Fourier coefficients of the real-valued version of $f$ over subsets of
coordinates of size greater than $30w\lceil\log_2(3/\varepsilon)\rceil$ is at most
$\varepsilon$.
\end{theorem}

\begin{theorem}[Low degree part eq sparse fourier approximation]
\lean{FABL.lowDegreePart_eq_sparseFourierApproximation}
\statementsource{Analysis-of-Boolean-Functions-by-Ryan-ODonnell}{pdf-ffbe9c7b3306-p028-b002}
\label{FABL.lowDegreePart_eq_sparseFourierApproximation}
\leanok
\uses{FABL.fourierCoeff, FABL.lowDegreePart, FABL.sparseFourierApproximation}
\difficulty{1}
% dataset-pipeline: stub -- prose pending build_informal_statements.py
Statement pending regeneration.
\end{theorem}

\begin{theorem}[Fourier coefficients of the low-degree part]
\lean{FABL.fourierCoeff_lowDegreePart}
\label{FABL.fourierCoeff_lowDegreePart}
\leanok
\uses{FABL.fourierCoeff, FABL.fourierCoeff_sparseFourierApproximation, FABL.lowDegreePart, FABL.lowDegreePart_eq_sparseFourierApproximation}
\difficulty{3}
Let $n$ and $k$ be nonnegative integers, let $f$ be a real-valued function on the
$n$-dimensional Boolean cube, and let $T$ be a finite set of its coordinates. The
Fourier coefficient at $T$ of the sum of the Fourier terms of $f$ supported on sets of
size at most $k$ equals the Fourier coefficient of $f$ at $T$ when $|T|\le k$, and
equals $0$ otherwise.
\end{theorem}

\begin{theorem}[Fourier norm of a low-degree part]
\lean{FABL.fourierOneNorm_lowDegreePart}
\label{FABL.fourierOneNorm_lowDegreePart}
\leanok
\uses{FABL.fourierCoeff, FABL.fourierCoeff_lowDegreePart, FABL.fourierOneNorm, FABL.lowDegreePart}
\difficulty{4}
Let $n$ be a nonnegative integer, $k$ a nonnegative integer, and $f$ a real-valued
function on the $n$-dimensional sign cube. The Fourier one-norm of the part of $f$
formed from Fourier terms indexed by subsets of size at most $k$ equals the sum of the
absolute values of the Fourier coefficients of $f$ over those subsets.
\end{theorem}

\begin{theorem}[Error after low-degree truncation]
\lean{FABL.uniformLpNorm_sub_lowDegreePart_sq}
\label{FABL.uniformLpNorm_sub_lowDegreePart_sq}
\leanok
\uses{FABL.fourierWeight, FABL.fourierWeightAbove, FABL.lowDegreePart, FABL.lowDegreePart_eq_sparseFourierApproximation, FABL.uniformLpNorm, FABL.uniformLpNorm_sub_sparseFourierApproximation_sq}
\difficulty{3}
Let $n,k\in\mathbb{N}$, and let $f:\{-1,1\}^n\to\mathbb{R}$. The square of the uniform
$2$-norm of the difference between $f$ and its degree-at-most-$k$ Fourier part equals
the sum of the squared Fourier coefficients of $f$ over all subsets of coordinates whose
sizes exceed $k$.
\end{theorem}

\begin{theorem}[Symmetry of squared difference norm]
\lean{FABL.uniformLpNorm_sub_sq_comm}
\label{FABL.uniformLpNorm_sub_sq_comm}
\leanok
\uses{FABL.uniformLpNorm, FABL.uniformLpNorm_two_sq_eq_expect_sq}
\difficulty{3}
Let $n\in\mathbb{N}$ and let $f,g:\{-1,1\}^n\to\mathbb{R}$. The squared uniform $2$-norm
of $f-g$ equals that of $g-f$.
\end{theorem}

\begin{definition}[Low-degree concentrating Fourier family]
\lean{FABL.lowDegreeL1ConcentratingFourierFamily}
\statementsource{Analysis-of-Boolean-Functions-by-Ryan-ODonnell}{pdf-ffbe9c7b3306-p102-b005}
\label{FABL.lowDegreeL1ConcentratingFourierFamily}
\leanok
\uses{FABL.l1ConcentratingFourierFamily, FABL.lowDegreePart}
Let $n$ be a natural number, let $f$ be a real-valued function on the $n$-dimensional
sign cube, let $k$ be a natural number, and let $\varepsilon$ be real. Define the
low-degree $L^1$-concentrating Fourier family of $f$ at degree $k$ and tolerance
$\varepsilon$ to be the $L^1$-concentrating Fourier family of the degree-$k$ low-degree
part of $f$, with tolerance $\varepsilon/4$.
\end{definition}

\begin{theorem}[Concentration on an $L^1$-selected family]
\lean{FABL.isFourierSpectrumConcentratedOn_lowDegreeL1ConcentratingFourierFamily}
\proofsource{Analysis-of-Boolean-Functions-by-Ryan-ODonnell}{
  pdf-ffbe9c7b3306-p102-b005,
  pdf-ffbe9c7b3306-p079-b003,
  pdf-ffbe9c7b3306-p069-b003
}
\label{FABL.isFourierSpectrumConcentratedOn_lowDegreeL1ConcentratingFourierFamily}
\leanok
\uses{FABL.IsFourierSpectrumConcentratedOn, FABL.IsFourierSpectrumConcentratedOn.transfer_of_uniformLpNorm_sub_sq_le, FABL.fourierWeightAbove, FABL.isFourierSpectrumConcentratedOn_l1ConcentratingFourierFamily, FABL.l1ConcentratingFourierFamily, FABL.lowDegreeL1ConcentratingFourierFamily, FABL.lowDegreePart, FABL.uniformLpNorm, FABL.uniformLpNorm_sub_lowDegreePart_sq, FABL.uniformLpNorm_sub_sq_comm}
\difficulty{4}
Let $n,k\ge 0$, let $f:\{\pm1\}^n\to\mathbb R$, and suppose $\varepsilon>0$. Assume the
total squared Fourier weight of $f$ on sets of cardinality greater than $k$ is at most
$\varepsilon/4$. Form the degree-at-most-$k$ Fourier truncation of $f$, and select the
sets $S$ for which $(\varepsilon/4)\le \|\widehat{g}\|_1|\widehat{g}(S)|$, where $g$ is
that truncation and $\|\widehat{g}\|_1$ is the sum of the absolute values of its Fourier
coefficients. Then the total squared Fourier weight of $f$ outside the selected family
is at most $\varepsilon$.
\end{theorem}

\begin{theorem}[Bound for large low degree Fourier coefficients]
\lean{FABL.card_lowDegreeL1ConcentratingFourierFamily_le}
\proofsource{Analysis-of-Boolean-Functions-by-Ryan-ODonnell}{pdf-ffbe9c7b3306-p087-b002}
\label{FABL.card_lowDegreeL1ConcentratingFourierFamily_le}
\leanok
\uses{FABL.card_l1ConcentratingFourierFamily_le, FABL.fourierOneNorm, FABL.fourierOneNorm_nonneg, FABL.lowDegreeL1ConcentratingFourierFamily, FABL.lowDegreePart}
\difficulty{4}
Let $n$ and $k$ be natural numbers and let $f$ be a real-valued function on
$\{−1,1\}^n$. For each subset $S$ of the coordinates, let $\widehat g(S)$ be the uniform
Fourier coefficient of the degree-at-most-$k$ part $g$ of $f$, and let
$\|g\|_1=\sum_S|\widehat g(S)|$. Given real numbers $\varepsilon,L$ with $\varepsilon>0$
and $L\ge0$, if $\|g\|_1\le L$, then the number of subsets $S$ satisfying
$\frac{\varepsilon}{4}\le \|g\|_1|\widehat g(S)|$ is at most $4L^2/\varepsilon$.
\end{theorem}

\begin{theorem}[Exists fourier concentrating family of low degree bounds]
\lean{FABL.exists_fourierConcentratingFamily_of_lowDegree_bounds}
\statementsource{Analysis-of-Boolean-Functions-by-Ryan-ODonnell}{
  pdf-ffbe9c7b3306-p087-b002,
  pdf-ffbe9c7b3306-p102-b005
}
\label{FABL.exists_fourierConcentratingFamily_of_lowDegree_bounds}
\leanok
\uses{FABL.IsFourierSpectrumConcentratedOn, FABL.IsFourierSpectrumConcentratedUpTo, FABL.card_lowDegreeL1ConcentratingFourierFamily_le, FABL.fourierCoeff, FABL.fourierOneNorm_lowDegreePart, FABL.fourierWeightAboveReal_natCast, FABL.isFourierSpectrumConcentratedOn_lowDegreeL1ConcentratingFourierFamily, FABL.lowDegreeL1ConcentratingFourierFamily}
\difficulty{3}
% dataset-pipeline: stub -- prose pending build_informal_statements.py
Statement pending regeneration.
\end{theorem}

\begin{definition}[DNF spectral concentration degree]
\lean{FABL.dnfSpectralConcentrationDegree}
\statementsource{Analysis-of-Boolean-Functions-by-Ryan-ODonnell}{
  pdf-ffbe9c7b3306-p101-b004,
  pdf-ffbe9c7b3306-p269-b005
}
\label{FABL.dnfSpectralConcentrationDegree}
\leanok
\uses{FABL.dnfSwitchingDepth}
For a natural number $w$ and a real number $\varepsilon$, define the DNF spectral
concentration degree to be $30w$ times the DNF switching depth at tolerance
$\varepsilon/4$.
\end{definition}

\begin{definition}[DNF spectral family size bound]
\lean{FABL.dnfSpectralFamilySizeBound}
\statementsource{Analysis-of-Boolean-Functions-by-Ryan-ODonnell}{pdf-ffbe9c7b3306-p102-b004}
\label{FABL.dnfSpectralFamilySizeBound}
\leanok
\uses{FABL.dnfSpectralConcentrationDegree}
For a natural number $w$ and a real number $\varepsilon$, define the DNF spectral family
size bound to be $\frac{4\bigl(2(20w)^{d(w,\varepsilon)}\bigr)^2}{\varepsilon}$, where
$d(w,\varepsilon)$ is the DNF spectral concentration degree at width $w$ and tolerance
$\varepsilon$.
\end{definition}

\begin{theorem}[Fourier concentration from switching bounds]
\lean{FABL.theorem4_25_of_switchingFailureProbability_bounds}
\proofsource{Analysis-of-Boolean-Functions-by-Ryan-ODonnell}{
  pdf-ffbe9c7b3306-p102-b004,
  pdf-ffbe9c7b3306-p102-b005
}
\label{FABL.theorem4_25_of_switchingFailureProbability_bounds}
\leanok
\uses{FABL.BooleanFunction, FABL.BooleanFunction.toReal, FABL.IsFourierSpectrumConcentratedOn, FABL.IsFourierSpectrumConcentratedUpTo, FABL.dnfSpectralConcentrationDegree, FABL.dnfSpectralFamilySizeBound, FABL.exists_fourierConcentratingFamily_of_lowDegree_bounds, FABL.fourierCoeff, FABL.switchingFailureProbability, FABL.theorem4_22_of_switchingFailureProbability_le_half_pow, FABL.theorem4_24_of_switchingFailureProbability_le_quarter}
\difficulty{4}
Let $n$ be a nonnegative integer and let $f:\{-1,1\}^n\to\{-1,1\}$. Let $w$ be a
positive integer and let $0<\varepsilon\le 1/2$. Suppose that for every integer $d\ge0$,
the probability that a random restriction with each coordinate left free independently
with probability $1/(10w)$ has decision-tree depth at least $d$ is at most $2^{-d}$, and
the corresponding probability for free-coordinate probability $1/(20w)$ is at most
$4^{-d}$. Then there is a finite family $\mathcal F$ of subsets of $[n]$ such that the
total squared Fourier coefficient mass outside $\mathcal F$ is at most $\varepsilon$,
and $|\mathcal F|\le 4\bigl(2(20w)^D\bigr)^2/\varepsilon$, where
$D=30w\left\lceil\log_2(12/\varepsilon)\right\rceil$ (with the ceiling taken as a
nonnegative integer).
\end{theorem}

\begin{theorem}[Håstad half switching bound]
\lean{FABL.hastadHalfSwitchingBound}
\proofsource{Analysis-of-Boolean-Functions-by-Ryan-ODonnell}{pdf-ffbe9c7b3306-p100-b005}
\label{FABL.hastadHalfSwitchingBound}
\leanok
\uses{FABL.BooleanFunction, FABL.HasCNFWidthLE, FABL.HasDNFWidthLE, FABL.hastadSwitchingLemma, FABL.switchingFailureProbability}
\difficulty{4}
Let $n$ be a natural number and let $f:\{-1,1\}^n\to\{-1,1\}$. Suppose there is a
positive integer $w$ such that $f$ has a disjunctive normal form or a conjunctive normal
form whose terms or clauses each have width at most $w$. For every natural number $d$,
the probability under the random restriction that leaves each coordinate free with
probability $1/(10w)$ that the restricted function has decision-tree depth at least $d$
is at most $(1/2)^d$.
\end{theorem}

\begin{theorem}[Quarter switching bound for formulas]
\lean{FABL.hastadQuarterSwitchingBoundForFormula}
\proofsource{Analysis-of-Boolean-Functions-by-Ryan-ODonnell}{pdf-ffbe9c7b3306-p100-b005}
\label{FABL.hastadQuarterSwitchingBoundForFormula}
\leanok
\uses{FABL.BooleanFunction, FABL.HasCNFWidthLE, FABL.HasDNFWidthLE, FABL.hastadSwitchingLemma, FABL.switchingFailureProbability}
\difficulty{4}
Let $n$ be a natural number and let $f:\{-1,1\}^n\to\{-1,1\}$. Suppose there is a
positive integer $w$ such that $f$ has a disjunctive normal form or a conjunctive normal
form whose terms or clauses each have width at most $w$. For every natural number $d$,
the probability under the random restriction that leaves each coordinate free with
probability $1/(20w)$ that the restricted function has decision-tree depth at least $d$
is at most $(1/4)^d$.
\end{theorem}

\begin{theorem}[Constant functions have no nonconstant Fourier weight]
\lean{FABL.isFourierSpectrumConcentratedUpTo_zero_of_constant}
\label{FABL.isFourierSpectrumConcentratedUpTo_zero_of_constant}
\leanok
\uses{FABL.BooleanFunction, FABL.BooleanFunction.toReal, FABL.IsFourierSpectrumConcentratedUpTo, FABL.expect_monomial, FABL.fourierCoeff, FABL.fourierWeight, FABL.fourierWeightAboveReal}
\difficulty{4}
Let $f$ be a Boolean function on the sign cube that is constant, and let $\varepsilon\ge
0$. The sum of squared Fourier coefficients indexed by nonempty subsets of coordinates
is at most $\varepsilon$.
\end{theorem}

\begin{theorem}[Fourier concentration for bounded-width DNF]
\lean{FABL.theorem4_22_of_pos}
\proofsource{Analysis-of-Boolean-Functions-by-Ryan-ODonnell}{
  pdf-ffbe9c7b3306-p101-b004,
  pdf-ffbe9c7b3306-p101-b005,
  pdf-ffbe9c7b3306-p101-b003,
  pdf-ffbe9c7b3306-p101-b002
}
\label{FABL.theorem4_22_of_pos}
\leanok
\uses{FABL.BooleanFunction, FABL.BooleanFunction.toReal, FABL.HasDNFWidthLE, FABL.IsFourierSpectrumConcentratedUpTo, FABL.dnfSwitchingDepth, FABL.hastadHalfSwitchingBound, FABL.theorem4_22_of_switchingFailureProbability_le_half_pow}
\difficulty{1}
Let $n$ be a nonnegative integer and let $f$ be a function from the $n$-dimensional sign
cube to the two signs. Suppose $f$ has a disjunctive normal form of width at most $w$,
where $w$ is a positive integer. For any real number $\varepsilon$ with
$0<\varepsilon\le 1$, the sum of squared Fourier coefficients of the real-valued version
of $f$ over subsets of coordinates of size greater than
$30w\lceil\log_2(3/\varepsilon)\rceil$ is at most $\varepsilon$.
\end{theorem}

\begin{theorem}[Theorem4 22]
\lean{FABL.theorem4_22}
\statementsource{Analysis-of-Boolean-Functions-by-Ryan-ODonnell}{
  pdf-ffbe9c7b3306-p101-b004,
  pdf-ffbe9c7b3306-p101-b005,
  pdf-ffbe9c7b3306-p101-b003
}
\label{FABL.theorem4_22}
\leanok
\uses{FABL.BooleanFunction, FABL.BooleanFunction.toReal, FABL.DNFFormula.toBooleanFunction, FABL.DNFFormula.toBooleanFunction_eq_of_width_eq_zero, FABL.DNFFormula.width, FABL.HasDNFWidthLE, FABL.IsFourierSpectrumConcentratedUpTo, FABL.dnfSwitchingDepth, FABL.isFourierSpectrumConcentratedUpTo_zero_of_constant, FABL.theorem4_22_of_pos}
\difficulty{4}
% dataset-pipeline: stub -- prose pending build_informal_statements.py
Statement pending regeneration.
\end{theorem}

\begin{theorem}[Low-degree Fourier coefficient bound for DNF]
\lean{FABL.theorem4_24}
\proofsource{Analysis-of-Boolean-Functions-by-Ryan-ODonnell}{
  pdf-ffbe9c7b3306-p102-b002,
  pdf-ffbe9c7b3306-p102-b003
}
\label{FABL.theorem4_24}
\leanok
\uses{FABL.BooleanFunction, FABL.BooleanFunction.toReal, FABL.HasDNFWidthLE, FABL.fourierCoeff, FABL.hastadQuarterSwitchingBoundForFormula, FABL.theorem4_24_of_switchingFailureProbability_le_quarter}
\difficulty{1}
Let $n$ be a nonnegative integer and let $f$ be a sign-valued Boolean function on the
$n$-dimensional sign cube. Suppose $f$ has a DNF representation of width at most the
positive integer $w$. For every nonnegative integer $k$, the sum of the absolute values
of the Fourier coefficients of the real-valued version of $f$ over all coordinate
subsets of size at most $k$ is at most $2(20w)^k$.
\end{theorem}

\begin{theorem}[Small Fourier family for bounded-width DNF]
\lean{FABL.theorem4_25}
\proofsource{Analysis-of-Boolean-Functions-by-Ryan-ODonnell}{
  pdf-ffbe9c7b3306-p102-b004,
  pdf-ffbe9c7b3306-p102-b005,
  pdf-ffbe9c7b3306-p102-b002,
  pdf-ffbe9c7b3306-p102-b003,
  pdf-ffbe9c7b3306-p101-b004,
  pdf-ffbe9c7b3306-p100-b005,
  pdf-ffbe9c7b3306-p087-b002,
  pdf-ffbe9c7b3306-p087-b003
}
\label{FABL.theorem4_25}
\leanok
\uses{FABL.BooleanFunction, FABL.BooleanFunction.toReal, FABL.HasDNFWidthLE, FABL.IsFourierSpectrumConcentratedOn, FABL.dnfSpectralFamilySizeBound, FABL.hastadHalfSwitchingBound, FABL.hastadQuarterSwitchingBoundForFormula, FABL.theorem4_25_of_switchingFailureProbability_bounds}
\difficulty{2}
Let $n$ be a nonnegative integer and let $f$ be a sign-valued Boolean function on the
$n$-dimensional sign cube. Suppose $f$ has a DNF representation of width at most an
integer $w\ge2$, and let $0<\varepsilon\le1/2$. Then there is a finite family $\mathcal
F$ of subsets of the coordinates such that the sum of squared Fourier coefficients of
the real-valued version of $f$ outside $\mathcal F$ is at most $\varepsilon$, and
\[|\mathcal F|\le \frac{4\left(2(20w)^{30w\lceil\log_2(12/\varepsilon)\rceil}\right)^2}{\varepsilon}.\]
\end{theorem}
